% ===========================================================================
%  sviluppatore/09-riferimento/02-costanti.tex
% ===========================================================================
\section{Le costanti e la configurazione}\label{sez:costanti}

\subsection{\file{config/quantsys.yaml}}

È il file centrale. Dopo una modifica serve riavviare la dashboard: non
viene riletto a caldo.

\begin{center}
\begin{tabular}{@{}L{46mm}R{18mm}L{72mm}@{}}
\toprule
\textbf{Chiave} & \textbf{Valore} & \textbf{Che cosa governa} \\
\midrule
\multicolumn{3}{@{}l}{\textit{\small \cod{data:}}}\\[2pt]
\rigaTenue
\cod{start\_date}        & \num{2010}-01-01 & Inizio del campione. \\
\cod{train\_years}       & \num{5}  & Ampiezza della finestra di
                                      addestramento. \\
\rigaTenue
\cod{test\_years}        & \num{1}  & Ampiezza della finestra di prova. \\
\cod{slide\_years}       & \num{1}  & Di quanto avanza la finestra. \\
\rigaTenue
\cod{warmup\_days}       & \num{300} & Barre saltate all'inizio. \\
\cod{min\_folds}         & \num{3}  & Minimo di finestre perché il
                                      walk-forward abbia senso. \\
\rigaTenue
\cod{benchmark\_symbol}  & SPY      & Il riferimento dei confronti. \\

\addlinespace
\multicolumn{3}{@{}l}{\textit{\small \cod{portfolio:}}}\\[2pt]
\rigaTenue
\cod{initial\_capital}   & \num{1000000} & Capitale iniziale. \\
\cod{portfolio\_target\_vol\_annual} & \num{0.12} & Volatilità annua
                                      obiettivo. \\
\rigaTenue
\cod{max\_position\_weight} & \num{0.05} & Tetto per strumento. \\
\cod{max\_asset\_class\_gross} & \num{1.00} & Tetto lordo per classe. \\
\rigaTenue
\cod{max\_gross\_leverage} & \num{1.0} & Tetto lordo totale. \\

\addlinespace
\multicolumn{3}{@{}l}{\textit{\small \cod{ensemble:}}}\\[2pt]
\rigaTenue
\cod{target\_vol}        & \num{0.10} & Volatilità a cui ogni stream viene
                                      portato \emph{prima} di essere pesato. \\
\cod{scale\_lookback}    & \num{252} & Barre della volatilità trailing. \\
\rigaTenue
\cod{regime\_window}     & \num{90}  & Finestra del regime. \\
\cod{regime\_sh\_low} / \cod{\_high} & \num{0.0} / \num{1.5} & Soglie di
                                      Sharpe per la modulazione. \\
\bottomrule
\end{tabular}
\end{center}

\begin{attenzione}{\cod{ensemble.target\_vol} non è un parametro di reporting}
Il commento nel file lo dice: «non è un dettaglio di reporting: lo stesso
fattore moltiplica il libro eseguito dal vivo». Uno stream a volatilità molto
bassa chiede molta leva per arrivare all'obiettivo, e il tetto resta
\cod{max\_gross\_leverage}. Cambiare questo valore cambia le dimensioni delle
posizioni su un conto reale (\vediSez{live}).
\end{attenzione}

\noindent
I valori della tabella sono quelli del file come si trova nel progetto; i
default del codice, usati quando una chiave manca, sono diversi e stanno in
\simbolo{AllocatorConfig} (\vediSez{allocatore}).

\subsection{Le costanti del motore}

\begin{center}
\begin{tabular}{@{}L{40mm}R{14mm}L{82mm}@{}}
\toprule
\textbf{Costante} & \textbf{Valore} & \textbf{Dove, e a che serve} \\
\midrule
\rigaTenue
\cod{SCALE\_MIN\_OBS} & \num{20} & \file{reporting/ensemble.py}. Sotto
  questa soglia non c'è stima di volatilità e il fattore vale \cod{1.0}. \\
\cod{SCALE\_MAX\_LEVA} & \num{10.0} & Idem. Rapporto massimo fra obiettivo e
  volatilità stimata: è il tetto alla leva quando la volatilità tende a
  zero. \\
\rigaTenue
\cod{max\_workers} & \num{1} & \file{execution/job\_manager.py}. I job sono
  serializzati (\vediSez{job}). \\
\cod{max\_log\_lines} & \num{5000} & Idem. Righe di registro conservate per
  ogni job. \\
\rigaTenue
\cod{min\_lookback\_days} & \num{252} & Default di
  \simbolo{StrategyDescriptor}. \\
\bottomrule
\end{tabular}
\end{center}

\subsection{Le variabili d'ambiente}

\begin{center}
\begin{tabular}{@{}L{38mm}L{100mm}@{}}
\toprule
\textbf{Variabile} & \textbf{Effetto} \\
\midrule
\rigaTenue
\cod{LYRA\_HOME} & Sposta la cartella di lavoro (\vediSez{dati}). \\
\cod{LYRA\_APP\_TOKEN} & Scavalca il token personale letto da file
  (\vediSez{ponte}). \\
\rigaTenue
\cod{FLASK\_HOST} & Default \cod{127.0.0.1}. Cambiarlo espone la dashboard
  (\vediSez{avvio}). \\
\cod{FLASK\_PORT} & Default \num{5000}. \\
\rigaTenue
\cod{GH\_REPO} & Repository per gli aggiornamenti automatici
  (\vediSez{rilascio}). \\
\bottomrule
\end{tabular}
\end{center}

\noindent
Le chiavi dei servizi esterni --- Polygon, FRED, Stripe, SMTP, Google,
Postgres --- \emph{non} stanno qui: le gestisce il sito e arrivano alla metà
locale attraverso il canale~\num{2} (\vediSez{due-meta}).
